% !TEX root = lectures.tex
\section{Lecture 2}
\subsection{Probability Basics}
We started class with a probability review. Most
was generic stuff. But, the following definition is a bit rare.
\begin{definition}
    $X_1, \dots, X_n$ are parwise ($k$-wise, respectively)
    independent of every pair ($k$-tuple, respectively) are independent (their probabilities multiply).
\end{definition}
We move the discussion to useful tail bounds to estimate tail probabilities.
\begin{theorem}[Markov's Bound]
    If $X$ is a nonnegative random variable,
    \[ \Pr{X \ge k \E{X}} \le 1/k \]
    If in addition, $X \le z$, then
    \[ \Pr{X \le c \E{X}} \le \frac{z - \E{X}}{z - c\E{X}} \]
\end{theorem}
\begin{theorem}[Chebyshev's Inequality]
    For any random variable,
    \[ \Pr{|X - \E{X}| \ge k \sigma} \le 1/k^2 \]
\end{theorem}
Recall that a Gaussian distribution $N(\mu, \sigma^2) = \frac{1}{\sqrt{2\pi\sigma^2}} e^{-\frac{(x - \mu)^2}{2 \sigma^2}}$
concentrates very well around its mean. In particular,
if $X \sim N(\mu, \sigma^2)$, $\Pr{|X - \mu| \ge t\sigma} \le O(e^{-t^2/2})$.
It turns out that there are other settings that mimic Gaussians (largely due to the central limit theorem),
\begin{theorem}[Chernoff/Hoeffding's Bound]
    (Bernstein's Inequality)  Let $X_1, \dots, X_n$ be independent and they are bounded (say, $X_i \in [-1, 1]$).
    Call $\E{X_i} = \mu_i$, $\sigma_i^2 = \Var X_i$ and $X = \sum_i X_i$. Suppose $k \le \frac{\sigma}{2}$.
    \[ \Pr{|X- \mu| \ge k \sigma} \le 2\exp\qty(-\frac{k^2}{4}) \]

    (Chernoff Bound) Let $X_1, \dots, X_n \in \{0, 1\}$ be independent and $p_i = \E{X_i}$. Then,
    \[ \Pr{\sum_i X_i \ge (1 + \epsilon) \mu} \le e^{-\frac{\epsilon^2 \mu}{3 + \epsilon}} \]
    \begin{proof}
        We will use the so-called moment method.
        Define $Y = e^{tX}$ and 
        \begin{align*}
            \E{Y} &= \E{e^{t \sum_i X_i}} = \E{\prod_i e^{tX_i}} = \prod_i \E{e^{tX_i}} \\
            &= \prod_i (p_i e^t + (1 - p_i)) = \prod_i (1 + p_i(e^t - 1)) \le \prod_i e^{p_i(e^t - 1)} = e^{\mu(e^t - 1)}
        \end{align*}
        Thus, by Markov's,
        \begin{align*}
            \Pr{X \ge (1 + \epsilon) \mu} = \Pr{Y \ge e^{t\mu(1 + \epsilon)}} \le \frac{e^{\mu (e^t - 1)}}{e^{\mu t(1 + \epsilon)}}
        \end{align*}
        Optimizing $t$, as $t = \log(1 + \epsilon)$ and using $\log (1+\epsilon) \approxeq \epsilon - \frac{\epsilon^2}{2}$,
        we get the bound.
    \end{proof}
\end{theorem}

\subsection{Hashing}
We have a universe $U$ and we wish to map it to a small set of size $n \ll |U|$ using
a function $h$. The data structures we want to support are insert($x$), delete($x$)
and search($x$), which have self explanatory behavior. We can use randomness to make our
data structure to be robust against adversarial behavior.

A simple idea is
to make $h$ a function a which maps each element as it arrives into
an independent random slot. If there are $m = O(n)$ elements
that come in, the load on any slot is $\le O(\log n)$ with high probability.
\begin{theorem}
    The load on any slot is at most constant with high probability for $m < n$.
\end{theorem}
We can start the proof as follows.
Let $L_x$ be the number of elements in the same slot as $x$.
Define $I_y$ be the indicator that $h(y) = h(x)$.
It's clear that $ L_x = 1 + \sum_{y \neq x} I_y$. Thus, $\E{L_x} = 1 + \sum_{y \neq x} \E{I_y} = 1 + \frac{m-1}{n} \le 2$
if $m - 1 \le n$. We will pause here to note that if we
reduce the amount of randomness, then we can get the same result.
Namely, let $\mathcal{H}$ be a family of hash functions
and $h$ is a uniformly random element of $\mathcal{H}$. Namely,
as long as the values of $h$ are pairwise independent, $\E{I_y} = \frac{1}{n}$
(in fact, this is an even weaker notion of pairwise independence).
